% AUTORIDAD DE FUENTE: Exposé IX, §7.4, líneas 2017--2080;
% índices de escaneo 420--423, folios 409--412. Las capas francesas D/C
% no presentan diferencias materiales.

\medskip
\noindent
7.4. Supongamos ahora dado un segundo esquema en grupos $A'$ sobre $S$
que satisface las mismas hipótesis que $A$, y supongamos que $A$ y $A'$
tienen fibras conexas. Sea $W$ una biextensión de $(A,A')$ por
$\mathbb G_{mS}$. Obtenemos así un emparejamiento (VIII 2.2)

\begin{equation}\tag{7.4.1}
T_{\ell}(A) \times T_{\ell}(A') \longrightarrow
T_{\ell}(\mathbb G_{mS}) = \mathbb Z_{\ell}(1)
\end{equation}
y, por restricción a las partes fijas, un emparejamiento

\begin{equation}\tag{7.4.2}
T_{\ell}(A)^{f} \times T_{\ell}(A')^{f}
 = T_{\ell}(\widehat A) \times T_{\ell}(\widehat {A'})
 \longrightarrow \mathbb Z_{\ell}(1) \qquad .
\end{equation}
Este también puede interpretarse como el emparejamiento definido por
la biextensión $\widehat W$ de $(\widehat A,\widehat {A'})$ por
$\widehat{\mathbb G}_{mS}$ obtenida de $W$ por completado (aquí, por
simplicidad, se define una biextensión de esquemas formales en grupos
como una familia coherente de biextensiones $W_n$ de $(A_n,A'_n)$ por
$\mathbb G_{mS_n}$). Ahora bien, sabemos (VIII 3.5) que $W_n$ proviene
de una biextensión $W_n^{ab}$ de $(B_n,B'_n)$, donde $B'_n$ es la parte
abeliana de $A'_n$, de modo que (7.4.2) se factoriza a través de

\begin{equation}\tag{7.4.3}
T_{\ell}(A)^{ab} \times T_{\ell}(A')^{ab}
 = T_{\ell}(\widehat B) \times T_{\ell}(\widehat {B'})
 \longrightarrow \mathbb Z_{\ell}(1) \qquad .
\end{equation}
Este es también el emparejamiento definido por la biextensión

\begin{equation*}
\widehat W^{ab}=(W_n^{ab})
\end{equation*}
de $(\widehat B,\widehat {B'})$ por $\widehat{\mathbb G}_{mS}$. Cuando
$\widehat B$ y $\widehat {B'}$ son algebraizables, de modo que se
dispone de los esquemas abelianos $B,B'$ sobre $S$, EGA \sisi{IV}{III} 5.1.4
demuestra que $\widehat W^{ab}$ proviene de una biextensión, definida
salvo isomorfismo único,

%%%%%%%%%%  scan idx 421  |  folio 410  %%%%%%%%%%

\begin{equation*}
W^{ab}\quad\text{de}\quad(B,B')\quad\text{por}\quad\mathbb G_{mS},
\end{equation*}
y el emparejamiento correspondiente

\begin{equation}\tag{7.4.4}
T_{\ell}(B) \times T_{\ell}(B')
 \longrightarrow T_{\ell}(\mathbb G_{mS})
\end{equation}
es canónicamente isomorfo a (7.4.3).

Denotando por $D(B')$ el esquema abeliano dual de $B'$, $W^{ab}$ define
un homomorfismo

\begin{equation}\tag{7.4.5}
u:B\longrightarrow D(B') \qquad ,
\end{equation}
cuyas propiedades reflejan las de los emparejamientos (7.4.4), es
decir, (7.4.3). Por ejemplo, el emparejamiento (7.4.3) es una dualidad
perfecta en un punto $t\in S$ si y solo si
$u_t:B_t\longrightarrow D(B')_t$ es una isogenia cuyo núcleo tiene
rango primo con $\ell$. Esto implica que el propio $u$ es una isogenia
(pues $S$ es local, y por tanto conexo), y $\ker u$ es plano y finito
sobre $S$, de modo que $\ker u$ tiene el mismo rango en todo punto. Así,
la condición que acabamos de considerar es independiente del punto
elegido $t$.

\medskip
\noindent
7.4.6. Esto termina, en particular, la demostración de 5.4. En efecto,
bajo las condiciones de ese teorema (donde se puede suponer que $S$ es
completo), la aplicación de la discusión precedente a $A^o$ y
${A'}^o$ demuestra que, para todo $\ell$, el emparejamiento (7.4.3) es
una dualidad perfecta en el punto genérico $\eta$, y por tanto también
en el punto $s$.

%%%%%%%%%%  scan idx 422  |  folio 411  %%%%%%%%%%

\medskip
\noindent
7.4.7. Las consideraciones precedentes aclaran asimismo el teorema de
ortogonalidad 5.2, sin recurrir al Lema \sisi{5.2.4}{5.2.5}. En efecto, en el curso de
7.4 hemos establecido que el emparejamiento $\sisi{(5.2.3.1)}{(5.2.4.2)}=(7.4.2)$ se
anula sobre las partes tóricas de ambos lados y, por tanto, se factoriza
a través del emparejamiento (7.4.3). Solo queda establecer que este
emparejamiento es no degenerado sobre las fibras genéricas de los
progrupos considerados, sabiendo que lo es sobre las fibras especiales.
Esto se deduce del hecho de que, si (7.4.5) induce una isogenia sobre
las fibras especiales, entonces es una isogenia.

\medskip
\noindent
7.4.8. Volvamos a las condiciones generales de 7.4, y supongamos dado
un abierto $U\subset S-S_o$ tal que $A_U$ y $A'_U$ sean
\underline{esquemas abelianos}, duales entre sí mediante la biextensión
$W_U$. Restringiendo (7.3.6), y la sucesión exacta análoga para $A'$, a
$U$, obtenemos filtraciones canónicas sobre los grupos de Barsotti--Tate
pro-$\ell$ $T_\ell(A_U)$ y $T_\ell(A'_U)$:

\begin{equation}\tag{7.4.9}
\left\{
\begin{aligned}
T_\ell(A_U)&\supset T_\ell(A_U)^f\supset T_\ell(A_U)^t\supset0,\\[4ex]
T_\ell(A'_U)&\supset T_\ell(A'_U)^f\supset T_\ell(A'_U)^t\supset0
\qquad .
\end{aligned}
\right.
\end{equation}

%%%%%%%%%%  scan idx 423  |  folio 412  %%%%%%%%%%

Afirmo que \underline{estas dos filtraciones son duales entre sí} con
respecto al emparejamiento (7.4.1), que aquí es una dualidad perfecta
(VIII 3.2); es decir, la relación (5.2.2) y su contraparte simétrica
siguen siendo válidas al reemplazar $A_K,A'_K$ por $A_U,A'_U$. Hemos
visto arriba que esta afirmación de ortogonalidad equivale asimismo a
decir que la biextensión $W^{ab}$ de $(B,B')$ por $\mathbb G_{mS}$ es
una dualidad perfecta sobre $U$ (o, equivalentemente, cuando $U$ es
denso en $S$, sobre $S$). Basta comprobar esta propiedad en todo punto
$\eta$ de $U$; el argumento habitual nos reduce entonces al caso ya
conocido en que $S$ es un rasgo de punto genérico $\eta$.

La relación de ortogonalidad también permite extender las fórmulas
(5.5.7), (5.5.8) y (5.5.9) al caso presente, reemplazando la restricción
al punto genérico $\eta$ por la restricción al abierto $U$. Por
consiguiente, 5.6 se extiende asimismo al caso presente. Nos abstenemos
de reescribir aquí, en este contexto, las fórmulas y el enunciado
precedentes.
